\chapter{Cadre du projet}
\markboth{Chapitre 1}{Cadre du projet}
\setcounter{chapter}{1}
\minitoc
\newpage

%% ═══════════════════════════════════════════════════════════════════════════════
\section{Introduction}

Ce premier chapitre établit le cadre général du projet. Nous présentons
l'entreprise d'accueil Solvinya Group, détaillons la problématique commerciale
qui justifie ce développement, exposons la solution proposée, et décrivons
la méthodologie de développement adoptée.

%% ═══════════════════════════════════════════════════════════════════════════════
\section{Présentation de l'organisme d'accueil}

\subsection{Solvinya Group}

Solvinya Group est une \textbf{ESN} (Entreprise de Services Numériques) opérant
principalement sur les marchés francophones~: France, Belgique, Luxembourg,
Suisse et Tunisie. L'entreprise propose des services de conseil, de développement
logiciel, d'externalisation IT et d'accompagnement en transformation numérique.

\begin{table}[H]
  \centering
  \begin{tabular}{L{4cm} L{9cm}}
    \toprule
    \rowcolor{lightBlue!50}
    \textbf{Critère} & \textbf{Description} \\
    \midrule
    Secteur & ESN — Entreprise de Services Numériques \\
    Marchés cibles & France, Belgique, Luxembourg, Suisse, Tunisie \\
    Services & Développement logiciel, Conseil IT, Externalisation \\
    Technologies & COBOL/Mainframe, SAP, Java, Python, DevOps, Cloud \\
    Clients types & Banques, Assurances, Groupes industriels, Secteur public \\
    \bottomrule
  \end{tabular}
  \caption{Profil de Solvinya Group}
  \label{tab:solvinya}
\end{table}

\subsection{Domaines d'activités}

Solvinya Group intervient principalement dans les domaines suivants~:

\begin{itemize}
  \item \textbf{Systèmes legacy~:} maintenance et migration de systèmes
        COBOL/Mainframe dans le secteur bancaire ;
  \item \textbf{Intégration ERP~:} déploiement et personnalisation de SAP ;
  \item \textbf{Développement applicatif~:} Java, Python, React ;
  \item \textbf{Data \& IA~:} projets Data Engineering, Machine Learning ;
  \item \textbf{DevOps \& Cloud~:} CI/CD, infrastructure AWS/Azure.
\end{itemize}

%% ═══════════════════════════════════════════════════════════════════════════════
\section{Présentation du projet}

\subsection{Problématique}

La croissance de Solvinya Group repose sur sa capacité à identifier et contacter
rapidement les entreprises qui ont un \textbf{besoin IT actif}. Ce besoin peut
prendre plusieurs formes (illustrées en figure~\vref{fig:signal_types})~:

\begin{figure}[H]
  \centering
  \begin{tikzpicture}[
    node distance=1.2cm,
    box/.style={draw=espritBlue!60, fill=lightBlue!40, rounded corners=6pt,
                minimum width=3.5cm, minimum height=1.2cm,
                font=\small\bfseries, align=center, inner sep=6pt},
    arrow/.style={-{Stealth[length=6pt]}, espritBlue!70, thick}
  ]
    %% Centre
    \node[draw=espritBlue, fill=espritBlue!10, rounded corners=8pt,
          minimum width=3cm, minimum height=1.5cm, font=\bfseries\color{espritBlue},
          align=center] (center) {Besoin IT\\actif};

    %% Nœuds autour
    \node[box, above=1.8cm of center] (ao)
      {Appel d'offres\\public / privé};
    \node[box, right=2.5cm of center] (rfp)
      {RFP / Demande de\\proposition (DDP)};
    \node[box, below=1.8cm of center] (sub)
      {Sous-traitance /\\Externalisation};
    \node[box, left=2.5cm of center] (part)
      {Partenariat /\\Co-traitance};
    \node[draw=warnColor!80, fill=warnColor!40, rounded corners=6pt,
          minimum width=3.5cm, minimum height=1.2cm, font=\small\bfseries,
          align=center, above right=1.2cm and 1.8cm of center] (rec)
      {Recrutement IT\\(signal indirect)};

    \draw[arrow] (center) -- (ao);
    \draw[arrow] (center) -- (rfp);
    \draw[arrow] (center) -- (sub);
    \draw[arrow] (center) -- (part);
    \draw[arrow, dashed] (center) -- (rec);
  \end{tikzpicture}
  \caption{Les 5 formes de signaux d'un besoin IT actif}
  \label{fig:signal_types}
\end{figure}

\vspace{0.3cm}

\textbf{État actuel (sans l'outil)~:} les commerciaux effectuent manuellement
cette veille, ce qui représente~:
\begin{itemize}
  \item 3 à 5 heures par semaine de recherche manuelle sur LinkedIn, BOAMP et
        journaux officiels ;
  \item Un processus non systématique, dépendant du réseau personnel ;
  \item Des leads de qualité variable, souvent périmés au moment de la prise de contact ;
  \item Absence totale de couverture des marchés belge, luxembourgeois et suisse.
\end{itemize}

\subsection{Étude de l'existant}

\begin{table}[H]
  \centering
  \begin{tabular}{L{3cm} L{3.5cm} L{3.5cm} L{3cm}}
    \toprule
    \rowcolor{lightBlue!50}
    \textbf{Solution} & \textbf{Points forts} & \textbf{Limites} & \textbf{Coût} \\
    \midrule
    LinkedIn Sales Navigator & Interface intuitive, ciblage & Manuel, $99$/mois, pas d'AO publics & Élevé \\
    \addlinespace
    Hunter.io & Enrichissement email & Pas de détection de signaux & Moyen \\
    \addlinespace
    Apollo.io & Base de données contacts & Données US, peu FR/TN & Élevé \\
    \addlinespace
    Veille manuelle BOAMP & Gratuit, officiel & Non automatisé, ~3h/semaine & Gratuit \\
    \addlinespace
    \textbf{LeadGen 360+} & \textbf{Multi-source, NLP, LLM, AO publics FR/EU} & En développement & \textbf{Open source} \\
    \bottomrule
  \end{tabular}
  \caption{Comparaison des solutions de génération de leads B2B}
  \label{tab:existant}
\end{table}

\subsection{Solution proposée}

\textbf{LeadGen Francophone 360+} est une plateforme entièrement automatisée qui~:

\begin{enumerate}
  \item Scrape \textbf{5 sources} en parallèle : LinkedIn (offres emploi + posts B2B),
        BOAMP (marchés publics France), TED EU (marchés européens FR/BE/LU/CH) et Malt.fr ;
  \item Qualifie chaque prospect avec un \textbf{score 0--100} basé sur des règles métier ;
  \item Classe le texte via un modèle \textbf{NLP} (secteur + priorité) ;
  \item Génère des \textbf{messages de prospection personnalisés} via LLM (Groq) ;
  \item Automatise la collecte toutes les 12h via \textbf{n8n} et expose tout via une
        \textbf{API REST FastAPI}.
\end{enumerate}

%% ═══════════════════════════════════════════════════════════════════════════════
\section{Méthodologie de travail}

\subsection{Choix de la méthodologie Scrum}

Compte tenu de la nature itérative du projet (chaque collecteur est indépendant,
les exigences évoluent selon les tests réels sur les APIs), nous avons adopté
une approche \textbf{Scrum} adaptée.

\subsection{Planification des sprints}

\begin{table}[H]
  \centering
  \begin{tabular}{C{1.5cm} L{4cm} L{5cm} C{2cm}}
    \toprule
    \rowcolor{lightBlue!50}
    \textbf{Sprint} & \textbf{Objectif} & \textbf{Livrables} & \textbf{Durée} \\
    \midrule
    S1 & Infrastructure & Docker, PostgreSQL, FastAPI squelette & 1 semaine \\
    \addlinespace
    S2 & Collecteurs publics & BOAMP + TED EU + pipeline ingest & 2 semaines \\
    \addlinespace
    S3 & LinkedIn & linkedin-api wrapper, hiring + B2B & 2 semaines \\
    \addlinespace
    S4 & NLP & Classifier secteur + priorité TF-IDF & 1 semaine \\
    \addlinespace
    S5 & LLM & Groq client + message generator & 1 semaine \\
    \addlinespace
    S6 & Automatisation & n8n workflows + API collect endpoints & 1 semaine \\
    \addlinespace
    S7 & Tests \& Résultats & 815 prospects validés, rapport & 1 semaine \\
    \bottomrule
  \end{tabular}
  \caption{Planification Scrum du projet}
  \label{tab:sprints}
\end{table}

%% ═══════════════════════════════════════════════════════════════════════════════
\section{Conclusion}

Ce chapitre a posé le contexte du projet~: Solvinya Group a besoin d'automatiser
sa veille commerciale B2B sur les marchés francophones. La solution LeadGen 360+
répond à cette problématique en combinant scraping multi-source, NLP et génération
de contenu par LLM. Le chapitre suivant détaille l'analyse des besoins et
l'architecture conceptuelle du système.
